\documentclass[a4paper,11pt]{article}
\usepackage[top=16mm,bottom=20mm,left=16mm,right=16mm]{geometry}
\usepackage{fontspec}
\usepackage{xeCJK}
\setmainfont{Times New Roman}
\setCJKmainfont{Yu Gothic}
\setmonofont{Consolas}
\setCJKmonofont{Yu Gothic}
\usepackage{graphicx}
\usepackage{booktabs}
\usepackage{longtable}
\usepackage{array}
\usepackage{tabularx}
\usepackage{amsmath}
\usepackage{enumitem}
\usepackage{xstring}
\usepackage{tikz}
\usetikzlibrary{arrows.meta,positioning,shapes.geometric,fit}
\usepackage{xcolor}
\usepackage[unicode]{hyperref}
\hypersetup{
  hidelinks,
  pdftitle={Solar package flow workbook},
  pdfauthor={Codex}
}
\setlength{\parskip}{0.45em}
\setlength{\parindent}{1em}
\renewcommand{\arraystretch}{1.2}
\newif\ifanswers
\answerstrue
\newcommand{\fillin}[2]{\ifanswers{\color{black}\textbf{#1}}\else\underline{\hspace{#2}}\fi}
\newcommand{\flowbox}[1]{\fcolorbox{black}{white}{\parbox{0.95\linewidth}{#1}}}
\newcommand{\stagebox}[1]{\begin{center}\fcolorbox{black}{white}{\parbox{0.82\linewidth}{\centering #1}}\end{center}}
% Keep source names readable and allow line breaks at underscores.
\newcommand{\sourcepath}[1]{%
  \begingroup
  \StrSubstitute{#1}{\_}{\allowbreak\_}[\cleanpath]%
  {\ttfamily\cleanpath}%
  \endgroup
}
\title{改良版ソーラーカーパッケージ\\処理フロー体験ワークシート}
\author{solar\_ws0129-main}
\date{2026-07-16}

\begin{document}
\maketitle

\section*{使い方}
本書は、現在の改良版ソーラーカーパッケージについて、
\begin{itemize}[leftmargin=1.5em]
  \item PowerShell で何を叩くと、
  \item どの launch / script / node が動き、
  \item どの topic / CSV / dashboard / planner 更新へ流れるか
\end{itemize}
を、添付の ``MPC 処理フロー体験'' と同じく、手を動かして追える形でまとめたワークシートである。

PASSO 系列は本書では扱わない。対象はソーラーカー系の各モード、
\sourcepath{sim}、\sourcepath{measure}、\sourcepath{live}、\sourcepath{live\_wifi}、
\sourcepath{forecast}、\sourcepath{identify} である。

\begin{itemize}[leftmargin=1.5em]
  \item デフォルトは回答表示ありである。
  \item 穴埋めだけにしたい場合は、プリアンブルの \verb|\answerstrue| を \verb|\answersfalse| に変更する。
  \item 本書は 2026-07-16 時点のワークツリー、\sourcepath{SolarSim.ps1}、install 側の
  \sourcepath{scripts/solar\_control.sh}、\sourcepath{export\_rqt\_graph.py}、および現在確認できる core node 実装を根拠にしている。
\end{itemize}

\newpage
\tableofcontents
\newpage

\section{まず最初に押さえるべき事実}
\subsection{入口は 1 つではなく 3 層ある}
現行の操作系は、大きく次の 3 層に分かれる。

\stagebox{Windows 側入口: \sourcepath{SolarSim.ps1}}
\begin{center}$\Downarrow$\end{center}
\stagebox{WSL 側ルータ: \sourcepath{scripts/solar\_control.sh}}
\begin{center}$\Downarrow$\end{center}
\stagebox{\sourcepath{ros2 launch} または \sourcepath{python scripts/...}}
\begin{center}$\Downarrow$\end{center}
\stagebox{ROS 2 node 群と成果物: planner / logger / dashboard / bridge / CSV / JSON}

\subsection{入口と実装本体は source 側で対応している}
\flowbox{
  \textbf{2026-07-16 の監査結果:} \sourcepath{SolarSim.ps1} と \sourcepath{scripts/solar\_control.sh} が使う
  \sourcepath{solar\_measurement.launch.py}、\sourcepath{solar\_race\_live.launch.py}、
  \sourcepath{solar\_race\_live\_wifi.launch.py} は、いずれも \sourcepath{launch/} に存在する。
  telemetry / command / wind / state bridge も \sourcepath{mpc\_solarcar/} に実装され、
  \sourcepath{setup.py} の entry point と一致する。本書では、
  \begin{itemize}[leftmargin=1.5em]
    \item \textbf{PowerShell 入口から launch まで}
    \item \textbf{各 node の topic、数式、保存成果物}
  \end{itemize}
  を同じ source tree に基づいて追跡する。
}

\subsection{本書の根拠ファイル}
\begin{longtable}{p{0.30\linewidth}p{0.64\linewidth}}
\toprule
根拠ファイル & 本書での役割 \\
\midrule
\endhead
\sourcepath{SolarSim.ps1} & Windows から叩く最上位入口。action / mode / profile を WSL に渡す。 \\
\sourcepath{scripts/solar\_control.sh} & ROS起動、sim、forecast、identify、fit、learn、auditを振り分ける本体。 \\
\sourcepath{launch/*.launch.py} & sim / measure / live / live\_wifi の node 構成を定義する直接根拠。 \\
\sourcepath{scripts/solar\_sim.py} & オフライン大会全体 simulation。CSV、detail CSV、upper plan CSV、HTML、JSON を出力。 \\
\sourcepath{project\_packages/bwsc2027\_template/profile.yaml} & 改良版 profile 構造の代表例。各 mode がどのパラメータ群を使うかの根拠。 \\
\sourcepath{mpc\_solarcar/mpc\_node.py} & 上位・下位 planner の中心。通常1時間の上位再計画と 1 Hz lower command の根拠。 \\
\sourcepath{mpc\_solarcar/dashboard\_node.py} & dashboard が何を受け取り、何を API として見せるかの根拠。 \\
\sourcepath{mpc\_solarcar/distance\_node.py} & \sourcepath{/vehicle/speed\_kmh} から \sourcepath{/vehicle/s\_km} を積分する根拠。 \\
\sourcepath{mpc\_solarcar/grade\_node.py} & GPS / altitude / distance から勾配を出す根拠。 \\
\sourcepath{mpc\_solarcar/solar\_logger\_node.py} & 太陽車ログ CSV の列と topic 対応の根拠。 \\
\sourcepath{mpc\_solarcar/solar\_preflight\_node.py} & solar telemetry / planner freshness から system state / diag / health を作る。 \\
\sourcepath{scripts/export\_rqt\_graph.py} & mode 別 core node 群の graph export。 \\
\sourcepath{scripts/fetch\_weather\_forecast.py} & forecast action が何をするかの根拠。 \\
\sourcepath{scripts/run\_identification\_pipeline.py} & identify action の軽量 pipeline の根拠。 \\
\sourcepath{scripts/run\_vehicle\_identification.py} & 改良版 generic replay MLE identification の根拠。 \\
\bottomrule
\end{longtable}

\section{全体フローの骨格}
\subsection{PowerShell から見た action と mode}
\begin{center}
\begin{tabularx}{\linewidth}{>{\raggedright\arraybackslash}p{0.16\linewidth} >{\raggedright\arraybackslash}p{0.16\linewidth} X}
\toprule
action & mode & 到達先 \\
\midrule
\sourcepath{up} & \sourcepath{sim / measure / live / live\_wifi} & \sourcepath{build -> stop -> start}。最終的に \sourcepath{ros2 launch mpc\_solarcar <launch\_file>}。 \\
\sourcepath{start} & 同上 & 既存 build を前提に、その mode の launch だけ起動。 \\
\sourcepath{graph} & 同上 & \sourcepath{export\_rqt\_graph.py --mode <mode>}。graph exporter が想定 node 群を待って画像化。 \\
\sourcepath{simulate} & \sourcepath{sim} を主対象 & \sourcepath{python scripts/solar\_sim.py --profile\_yaml ...}。ROS2 launch ではなくオフライン実行。 \\
\sourcepath{forecast} & 主に \sourcepath{live / live\_wifi} 前準備 & \sourcepath{python fetch\_weather\_forecast.py --profile\_yaml ...}。Open-Meteo 取得。 \\
\sourcepath{identify} & profile ベース & \sourcepath{python run\_identification\_pipeline.py --profile\_yaml ...}。軽量な同定 pipeline。 \\
\sourcepath{fit} & profile ベース & 実ログを用いる高品質 segmented MLE と map shape fit。 \\
\sourcepath{learn} & profile ベース & CEMで上位目的重みを探索し、検証profileとTensorBoard logを保存。 \\
\sourcepath{audit} & 全体 & inventory、static contract audit、pytestを連続実行。 \\
\sourcepath{package / blank} & Windows直接 & 検証済みsolar-only版 / 同機能の空版を再生成。 \\
\sourcepath{log} & 任意 & \sourcepath{.run/solar\_<mode>.log} を tail。 \\
\bottomrule
\end{tabularx}
\end{center}

\subsection{WSL 側ルータでの mode -> launch 対応}
\begin{center}
\begin{tabular}{ll}
\toprule
mode & launch file（\sourcepath{solar\_control.sh} が期待する名前） \\
\midrule
\sourcepath{sim} & \sourcepath{solarcar\_sim.launch.py} \\
\sourcepath{measure} & \sourcepath{solar\_measurement.launch.py} \\
\sourcepath{live} & \sourcepath{solar\_race\_live.launch.py} \\
\sourcepath{live\_wifi} & \sourcepath{solar\_race\_live\_wifi.launch.py} \\
\bottomrule
\end{tabular}
\end{center}

\subsection{graph exporter が mode ごとに期待する core node}
\begin{center}
\begin{tabularx}{\linewidth}{>{\raggedright\arraybackslash}p{0.14\linewidth} X}
\toprule
mode & 期待 core node \\
\midrule
\sourcepath{sim} & \sourcepath{gps\_sim\_node}, \sourcepath{mpc\_node}, \sourcepath{solar\_state\_node}, \sourcepath{dashboard\_node} \\
\sourcepath{measure} & \sourcepath{dashboard\_node}, \sourcepath{solar\_logger\_node}, \sourcepath{distance\_node}, \sourcepath{grade\_node} \\
\sourcepath{live} & \sourcepath{mpc\_node}, \sourcepath{dashboard\_node}, \sourcepath{solar\_logger\_node}, \sourcepath{speed\_command\_bridge\_node} \\
\sourcepath{live\_wifi} & \sourcepath{mpc\_node}, \sourcepath{dashboard\_node}, \sourcepath{solar\_logger\_node}, \sourcepath{speed\_command\_bridge\_node}, \sourcepath{telemetry\_text\_bridge\_node}, \sourcepath{wind\_correction\_node} \\
\bottomrule
\end{tabularx}
\end{center}

\subsection{問題 1: 入口の穴埋め}
\begin{enumerate}[leftmargin=1.5em]
  \item Windows 側で最初に叩く入口は \fillin{\sourcepath{SolarSim.ps1}}{32mm} である。
  \item WSL 側で action / mode を解釈する本体は \fillin{\sourcepath{scripts/solar\_control.sh}}{48mm} である。
  \item \sourcepath{simulate} action は \fillin{\sourcepath{ros2 launch}}{22mm} ではなく、\fillin{\sourcepath{python scripts/solar\_sim.py}}{42mm} を直接呼ぶ。
  \item \sourcepath{live\_wifi} の graph exporter が追加で期待する bridge 系 node は
  \fillin{\sourcepath{telemetry\_text\_bridge\_node}}{44mm} と
  \fillin{\sourcepath{wind\_correction\_node}}{34mm} である。
\end{enumerate}

\section{profile.yaml の読み方}
\subsection{改良版 profile のトップレベル構造}
改良版 profile の代表例 \sourcepath{project\_packages/bwsc2027\_template/profile.yaml} では、
次のブロックが入口になる。

\begin{center}
\begin{tabularx}{\linewidth}{>{\raggedright\arraybackslash}p{0.18\linewidth} X}
\toprule
ブロック & 主な役割 \\
\midrule
\sourcepath{meta} & package 名、目的、注意事項。 \\
\sourcepath{paths} & route / weather / schedule / maps の実ファイル。多くの script がここを入口にする。 \\
\sourcepath{runtime} & dashboard host/port、forecast 時間解釈の基準。 \\
\sourcepath{logging} & 共通 log 置き場と rate。 \\
\sourcepath{simulation} & start 時刻、初期 SoC、出力先、prerace 実行条件。 \\
\sourcepath{measurement} & \sourcepath{distance\_node} と \sourcepath{grade\_node} の使い方。 \\
\sourcepath{identification} & 同定 input / output ディレクトリ。 \\
\sourcepath{live} & 実運用用。weather / autocal / command bridge / WiFi / wind model を束ねる。 \\
\sourcepath{model} & 車体質量、CdA、Crr、PV、battery 定数。 \\
\sourcepath{mpc} & 上位 / 下位 MPC の地平、重み、制約、再計画周期。 \\
\bottomrule
\end{tabularx}
\end{center}

\subsection{mode ごとに読む profile ブロック}
\begin{center}
\begin{tabularx}{\linewidth}{>{\raggedright\arraybackslash}p{0.14\linewidth} X}
\toprule
mode & 強く効くブロック \\
\midrule
\sourcepath{sim} & \sourcepath{paths}、\sourcepath{simulation}、\sourcepath{model}、\sourcepath{mpc}、\sourcepath{runtime} \\
\sourcepath{measure} & \sourcepath{measurement}、\sourcepath{logging}、必要なら \sourcepath{runtime} \\
\sourcepath{live} & \sourcepath{live.weather}、\sourcepath{live.autocal}、\sourcepath{live.command\_bridge}、\sourcepath{model}、\sourcepath{mpc} \\
\sourcepath{live\_wifi} & \sourcepath{live.wifi\_bridge}、\sourcepath{live.wind\_model}、\sourcepath{live.command\_bridge}、\sourcepath{live.weather}、\sourcepath{model}、\sourcepath{mpc} \\
\sourcepath{forecast} & \sourcepath{paths.route\_waypoints\_csv}、\sourcepath{paths.forecast\_csv}、\sourcepath{live.weather}、\sourcepath{runtime} \\
\sourcepath{identify} & \sourcepath{identification}、\sourcepath{paths}、\sourcepath{model} \\
\bottomrule
\end{tabularx}
\end{center}

\subsection{問題 2: profile の対応づけ}
\begin{enumerate}[leftmargin=1.5em]
  \item dashboard の host / port は主に \fillin{\sourcepath{runtime}}{18mm} ブロックで決まる。
  \item offline sim の初期 SoC、開始時刻、出力先は主に \fillin{\sourcepath{simulation}}{20mm} ブロックで決まる。
  \item 実運用の command smoothing や accel 制限は
  \fillin{\sourcepath{live.command\_bridge}}{38mm} に置く。
  \item WiFi 受送信先 IP / port は \fillin{\sourcepath{live.wifi\_bridge}}{34mm} に置く。
\end{enumerate}

\newpage
\section{sim モード}
\subsection{sim モードには 2 種類ある}
\begin{enumerate}[leftmargin=1.5em]
  \item \textbf{launch 型 sim}: \sourcepath{SolarSim.ps1 -Mode sim -Action up} で起動する ROS2 runtime。
  \item \textbf{offline 型 sim}: \sourcepath{SolarSim.ps1 -Mode sim -Action simulate} で実行する
  \sourcepath{scripts/solar\_sim.py}。
\end{enumerate}

この 2 つは同じ profile を参照するが、役割が違う。
\begin{itemize}[leftmargin=1.5em]
  \item launch 型 sim は topic を流し、planner と dashboard を起動する。
  \item offline 型 sim は大会全体 replay / planning 用の CSV、HTML、JSON を出力する。
\end{itemize}

\subsection{launch 型 sim の流れ}
\stagebox{\sourcepath{SolarSim.ps1 -Action up -Mode sim}
$\rightarrow$ \sourcepath{solarcar\_sim.launch.py}}
\begin{center}$\Downarrow$\end{center}
\stagebox{\sourcepath{gps\_sim\_node} / \sourcepath{solar\_state\_node}
$\leftrightarrow$ \sourcepath{mpc\_node}（次周期へ speed command feedback）}
\begin{center}$\Downarrow$\end{center}
\stagebox{\sourcepath{/sim/gps}、\sourcepath{/vehicle/*}、\sourcepath{/planner/*}
$\rightarrow$ \sourcepath{dashboard\_node}}

\subsection{launch 型 sim で直接確認できる node}
\sourcepath{launch/solarcar\_sim.launch.py} で確認できるのは次の 3 node である。
\begin{enumerate}[leftmargin=1.5em]
  \item \sourcepath{gps\_sim\_node}:
  \sourcepath{inputs/route\_waypoints.csv} と \sourcepath{/planner/speed\_cmd} を使って \sourcepath{/sim/gps} を publish する。
  \item \sourcepath{mpc\_node}:
  forecast CSV、route profile、speed profile、drive schedule、maps、params を読み、
  \sourcepath{/planner/speed\_cmd}、\sourcepath{/planner/upper\_plan}、
  \sourcepath{/planner/lower\_plan}、\sourcepath{/planner/status} などを publish する。
  \item \sourcepath{dashboard\_node}: \sourcepath{/api/state} を持つ HTTP サーバで、planner と vehicle の主要状態を可視化する。
\end{enumerate}

\subsection{offline 型 sim の流れ}
\sourcepath{scripts/solar\_sim.py} は ROS2 launch を通らず、profile を読んで大会全体をロールアウトする。

\stagebox{\sourcepath{profile.yaml}: route / weather / maps / schedule / model / MPC}
\begin{center}$\Downarrow$\end{center}
\stagebox{\sourcepath{scripts/solar\_sim.py}: full-course optimization と execution replay}
\begin{center}$\Downarrow$\end{center}
\stagebox{summary CSV / detail CSV / upper plan CSV / HTML / manifest JSON / resolved YAML}

\subsection{offline 型 sim で出る主成果物}
\begin{longtable}{p{0.28\linewidth}p{0.65\linewidth}}
\toprule
成果物 & 内容 \\
\midrule
\endhead
\sourcepath{out\_csv} & 時刻、距離、speed command、実行速度、SoC、battery 温度、G\_poa、外気温、向かい風。 \\
\sourcepath{out\_detail\_csv} & 電流、電圧、PV 出力、P\_pack、各抵抗力、drive / regen 効率など、再現確認用の詳細系列。 \\
\sourcepath{out\_plan\_csv} & 上位プランナが各再計画時刻で採用した segment 列。 \\
\sourcepath{report\_html} & 人間がブラウザで見る概要レポート。 \\
\sourcepath{summary\_json} & finish 到達、final distance、final SoC、CPU 時間などの要約。 \\
\sourcepath{resolved\_yaml} & override を織り込んだ実行時 profile の写し。 \\
\sourcepath{latest\_simulation\_run.json} & 最新 run の所在を追跡する manifest。 \\
\bottomrule
\end{longtable}

\subsection{sim で最重要な topic}
\begin{center}
\begin{tabularx}{\linewidth}{>{\raggedright\arraybackslash}p{0.26\linewidth} >{\raggedright\arraybackslash}p{0.17\linewidth} X}
\toprule
topic & 主 publisher & 意味 \\
\midrule
\sourcepath{/planner/speed\_cmd} & \sourcepath{mpc\_node} & planner が今出した速度指令。 \\
\sourcepath{/planner/upper\_speed\_cmd} & \sourcepath{mpc\_node} & 上位 planner 出力の代表値。 \\
\sourcepath{/planner/upper\_plan} & \sourcepath{mpc\_node} & 上位 plan の配列。dashboard に渡る。 \\
\sourcepath{/planner/lower\_plan} & \sourcepath{mpc\_node} & 下位 plan の配列。 \\
\sourcepath{/planner/status} & \sourcepath{mpc\_node} & 状態の数値サマリ。 \\
\sourcepath{/planner/env} & \sourcepath{mpc\_node} & planner が見ている環境量。 \\
\sourcepath{/planner/metrics} & \sourcepath{mpc\_node} & 電力・制約関連メトリクス。 \\
\sourcepath{/sim/gps} & \sourcepath{gps\_sim\_node} & route 上を進む模擬 GPS。 \\
\bottomrule
\end{tabularx}
\end{center}

\subsection{問題 3: sim フロー追跡}
\begin{enumerate}[leftmargin=1.5em]
  \item launch 型 sim で \sourcepath{/planner/speed\_cmd} を受けて route 上を進める node は
  \fillin{\sourcepath{gps\_sim\_node}}{28mm} である。
  \item offline 型 sim で大会全体の詳細な P\_pack、I、V まで出るのは
  \fillin{\sourcepath{out\_detail\_csv}}{26mm} である。
  \item 通常1時間の上位再計画と 1 Hz lower command の本体は \fillin{\sourcepath{mpc\_node}}{22mm} にある。
  \item dashboard の HTTP API endpoint は \fillin{\sourcepath{/api/state}}{20mm} である。
\end{enumerate}

\section{measure モード}
\subsection{measure モードの位置づけ}
measure は、実車や試験系から来る測定を
\begin{itemize}[leftmargin=1.5em]
  \item 距離に積分し、
  \item 勾配に変換し、
  \item ログに残し、
  \item dashboard で見えるようにする
\end{itemize}
ための mode である。

\flowbox{
  \textbf{実装根拠:} measurement launch が
  \sourcepath{dashboard\_node}、\sourcepath{solar\_logger\_node}、
  \sourcepath{distance\_node}、\sourcepath{grade\_node} を
  profile の \sourcepath{measurement} ブロックで起動する。
  本節は推測ではなく、launch と各 node source を直接追跡した処理である。
}

\subsection{measure の想定 runtime}
\stagebox{実車計測: \sourcepath{/vehicle/speed\_kmh} / GPS / altitude / battery}
\begin{center}$\Downarrow$\end{center}
\stagebox{\sourcepath{distance\_node}: \sourcepath{/vehicle/s\_km}
\quad+\quad \sourcepath{grade\_node}: \sourcepath{/vehicle/grade}}
\begin{center}$\Downarrow$\end{center}
\stagebox{\sourcepath{solar\_logger\_node}: CSV 記録
\quad+\quad \sourcepath{dashboard\_node}: 一画面表示}

\subsection{distance\_node がしていること}
\sourcepath{mpc\_solarcar/distance\_node.py} は、\sourcepath{/vehicle/speed\_kmh} を受け、
台形則に近い形で
\[
  s_{k+1} = s_k + \frac{v_k + v_{k+1}}{2}\frac{\Delta t}{3600}
\]
を積分し、\sourcepath{/vehicle/s\_km} を publish する。

したがって measure や live で距離が必要なら、\sourcepath{/vehicle/s\_km} を外から直接入れるか、
あるいは \sourcepath{distance\_node} を必ず併走させる。

\subsection{grade\_node がしていること}
\sourcepath{mpc\_solarcar/grade\_node.py} は、
\begin{itemize}[leftmargin=1.5em]
  \item \sourcepath{/vehicle/s\_km}
  \item \sourcepath{/vehicle/speed\_kmh}
  \item \sourcepath{/vehicle/altitude\_m} または GPS altitude
\end{itemize}
を受ける。標高は EMA で平滑化し、距離差が十分たまったときに勾配を更新して
\sourcepath{/vehicle/grade} を出す。

\subsection{solar\_logger\_node が残す列}
\sourcepath{solar\_logger\_node} は、少なくとも次を記録する。
\begin{itemize}[leftmargin=1.5em]
  \item ROS 時刻、UTC 壁時計、車速、距離、標高、勾配
  \item battery SoC / 温度 / 電流 / 電圧、PV 電力
  \item 実測・補正向かい風、風速・風向・進行方位、planner environment
  \item 上位 / 下位 speed command、drive mode、model metrics、online 校正値
  \item WiFi solar / chase 生文字列、system health / state / diag、各 node status
\end{itemize}

\subsection{問題 4: measure の要点}
\begin{enumerate}[leftmargin=1.5em]
  \item \sourcepath{/vehicle/s\_km} を生成する node は \fillin{\sourcepath{distance\_node}}{28mm} である。
  \item \sourcepath{/vehicle/grade} を出す node は \fillin{\sourcepath{grade\_node}}{24mm} である。
  \item CSV を残す node は \fillin{\sourcepath{solar\_logger\_node}}{30mm} である。
  \item measure mode の expected core node は
  \\[1mm]\fillin{\sourcepath{dashboard\_node, solar\_logger\_node, distance\_node, grade\_node}}{80mm} である。
\end{enumerate}

\section{live モード}
\subsection{live モードの位置づけ}
live は、実走中に planner を回し続け、測定を受けて速度指令を出す mode である。
改良版 profile では \sourcepath{live} ブロックの下に、
\begin{itemize}[leftmargin=1.5em]
  \item \sourcepath{weather}
  \item \sourcepath{autocal}
  \item \sourcepath{command\_bridge}
  \item \sourcepath{wind\_model}
\end{itemize}
がある。

\flowbox{
  \textbf{実装根拠:} live launch は
  weather、MPC、自動校正、指令 bridge、logger、dashboard を profile から構成する。
  \sourcepath{mpc\_node.py} の publish / subscribe と合わせ、
  実走時の経路を source で確認できる。
}

\subsection{live の想定流れ}
\stagebox{実測 topic: speed / distance / SoC / temperature / current / voltage}
\begin{center}$\Downarrow$\end{center}
\stagebox{\sourcepath{mpc\_node}: remaining-race upper MPC と 1 Hz lower MPC}
\begin{center}$\Downarrow$\end{center}
\stagebox{\sourcepath{speed\_command\_bridge\_node}: timeout / filter / rate limit / quantize}
\begin{center}$\Downarrow$\end{center}
\stagebox{実車 command。logger と dashboard は各段の同一 timestamp topic を並行購読}

\subsection{mpc\_node が live で受ける実測}
\sourcepath{mpc\_node} の solar mode では、少なくとも次の実測を subscribe する。
\begin{itemize}[leftmargin=1.5em]
  \item \sourcepath{/vehicle/s\_km}
  \item \sourcepath{/vehicle/speed\_kmh}
  \item \sourcepath{/vehicle/batt\_soc}
  \item \sourcepath{/vehicle/batt\_temp\_c}
  \item \sourcepath{/vehicle/batt\_current\_a}
  \item \sourcepath{/vehicle/batt\_voltage\_v}
\end{itemize}
上位 timer は 1 Hz で再計画条件を監視するが、重い上位最適化は
\sourcepath{upper\_replan\_sec}（通常3600秒）を満たした時に行う。
階層型が有効なら lower timer が \sourcepath{lower\_rate\_hz}（通常1 Hz）で指令を更新する。

\subsection{command bridge に期待される役割}
profile の \sourcepath{live.command\_bridge} には、
\begin{itemize}[leftmargin=1.5em]
  \item \sourcepath{publish\_rate\_hz}
  \item \sourcepath{input\_timeout\_sec}
  \item \sourcepath{safe\_speed\_kmh}
  \item \sourcepath{startup\_hold\_sec}
  \item \sourcepath{filter\_tau\_sec}
  \item \sourcepath{accel\_limit\_kmhps}
  \item \sourcepath{decel\_limit\_kmhps}
  \item \sourcepath{speed\_deadband\_kmh}
  \item \sourcepath{speed\_quantize\_step\_kmh}
  \item \sourcepath{max\_output\_speed\_kmh}
  \item \sourcepath{drive\_mode\_min\_hold\_sec}
\end{itemize}
があり、実運用時に ``出力が暴れない'' ようにする層である。

\subsection{weather / autocal / wind model の役割}
\begin{itemize}[leftmargin=1.5em]
  \item \sourcepath{live.weather}: 実走中の位置または fallback 緯度経度を使い、一定周期で予報を取得する想定。
  \item \sourcepath{live.autocal}: 走行中 / 停止中 / 昼夜条件に応じて、solar gain、drive power gain、補機電力推定をじわっと更新する想定。
  \item \sourcepath{live.wind\_model}: 風予報の分散成長、相関距離、planning quantile を使い、補正 forecast を planner に供給する想定。
\end{itemize}

\subsection{問題 5: live の穴埋め}
\begin{enumerate}[leftmargin=1.5em]
  \item live で \sourcepath{mpc\_node} が通常 1 Hz で指令更新するのは \fillin{下位 planner}{24mm} である。
  \item 実測の SoC は \fillin{\sourcepath{/vehicle/batt\_soc}}{30mm} topic で入る。
  \item 速度指令の平滑化や accel 制限は主に \fillin{\sourcepath{live.command\_bridge}}{38mm} が担う。
  \item live モードの graph exporter 期待 core node には
  \fillin{\sourcepath{speed\_command\_bridge\_node}}{44mm} が含まれる。
\end{enumerate}

\newpage
\section{live\_wifi モード}
\subsection{live\_wifi は live に WiFi 受送信層を足したもの}
live\_wifi では、live の planner / logger / dashboard に加え、
\begin{itemize}[leftmargin=1.5em]
  \item text telemetry bridge
  \item outbound command bridge
  \item wind correction
\end{itemize}
が追加される想定である。

\subsection{profile 上での主役}
\begin{center}
\begin{tabularx}{\linewidth}{>{\raggedright\arraybackslash}p{0.24\linewidth} X}
\toprule
ブロック & 役割 \\
\midrule
\sourcepath{live.wifi\_bridge} & 受信 bind host/port、送信先 host/port、送受信有効化、speed / distance / battery / wind のフィルタ時定数。 \\
\sourcepath{live.command\_bridge} & planner 出力を実車へ渡す直前の平滑化・安全化。 \\
\sourcepath{live.weather} & chase 側 GPS などから予報取得位置を決める。 \\
\sourcepath{live.wind\_model} & corrected forecast CSV を作る想定。 \\
\bottomrule
\end{tabularx}
\end{center}

\subsection{live\_wifi の設計フロー}
\stagebox{solar car + chase car: UTC timestamp 付き UTF-8 UDP JSON}
\begin{center}$\Downarrow$\end{center}
\stagebox{\sourcepath{telemetry\_text\_bridge\_node}: parse / UTC age-order gate / range gate / filter / ROS topic}
\begin{center}$\Downarrow$\end{center}
\stagebox{distance / grade / weather / wind correction
$\rightarrow$ \sourcepath{mpc\_node}}
\begin{center}$\Downarrow$\end{center}
\stagebox{command bridge $\rightarrow$ UDP command。
dashboard / logger は inbound、planner、outbound を並行記録}

\subsection{WiFi bridge の source と責務}
\flowbox{
  \textbf{実装済み:}
  \sourcepath{telemetry\_text\_bridge\_node.py}、
  \sourcepath{speed\_command\_bridge\_node.py}、
  \sourcepath{wind\_correction\_node.py}、
  \sourcepath{solar\_state\_node.py} は \sourcepath{mpc\_solarcar/} に存在する。
  \sourcepath{launch/solar\_race\_live\_wifi.launch.py} と \sourcepath{setup.py} の entry point も一致する。
  WiFi UTF-8 文字列受信から検証・filter、ROS topic、planner、指令平滑化、UDP 送信まで source で追跡できる。
  送信側 Raspberry Pi は NTP / chrony で UTC 同期し、毎 packet に
  \sourcepath{timestamp\_utc} または \sourcepath{ts\_unix} を付ける。
  既定では 5 秒超の遅延、2 秒超の未来時刻、重複・順序逆転を topic 更新前に拒否する。
}

\subsection{live\_wifi で追うべき代表 topic / データ}
\begin{itemize}[leftmargin=1.5em]
  \item inbound 側: speed、distance、battery、風、chase GPS
  \item 時刻監査: \sourcepath{/telemetry/solar\_source\_ts\_unix}、\sourcepath{/telemetry/chase\_source\_ts\_unix}、受信 UTC
  \item planner 内部: \sourcepath{/planner/speed\_cmd}、\sourcepath{/planner/upper\_plan}、\sourcepath{/planner/status}
  \item outbound 側: \sourcepath{output\_speed\_topic}、\sourcepath{output\_drive\_mode\_topic} または UDP / 文字列
  \item artifact 側: \sourcepath{live\_forecast\_raw.csv}、\sourcepath{live\_forecast\_corrected.csv}、runtime logs
\end{itemize}

\subsection{問題 6: live\_wifi の要点}
\begin{enumerate}[leftmargin=1.5em]
  \item WiFi の bind host / port を置くのは \fillin{\sourcepath{live.wifi\_bridge}}{34mm} である。
  \item corrected forecast CSV の出力先を持つのは
  \fillin{\sourcepath{live.wind\_model}}{30mm} である。
  \item live\_wifi の graph exporter 期待 core node は live に加えて、次の2 nodeを持つ。\\
  \fillin{telemetry\_text\_bridge\_node と wind\_correction\_node}{68mm}
\end{enumerate}

\section{forecast action}
\subsection{forecast は launch ではなく utility 実行}
\sourcepath{SolarSim.ps1 -Action forecast} は、最終的に
\sourcepath{scripts/fetch\_weather\_forecast.py --profile\_yaml ...}
を呼ぶ。

\subsection{何をしているか}
\begin{enumerate}[leftmargin=1.5em]
  \item profile を開く。
  \item \sourcepath{live.weather} と \sourcepath{runtime} から、forecast days、step minutes、timezone、tcell gain を決める。
  \item route waypoint CSV があれば先頭点の緯度経度を使う。
  \item 無ければ fallback 緯度経度を使う。
  \item Open-Meteo forecast を取得する。
  \item \sourcepath{paths.forecast\_csv} へ保存する。
\end{enumerate}

\subsection{forecast で触る代表パラメータ}
\begin{itemize}[leftmargin=1.5em]
  \item \sourcepath{forecast\_days}
  \item \sourcepath{step\_minutes}
  \item \sourcepath{timezone\_name}
  \item \sourcepath{fallback\_latitude / fallback\_longitude}
  \item \sourcepath{tcell\_gain}
\end{itemize}

\subsection{問題 7: forecast の穴埋め}
\begin{enumerate}[leftmargin=1.5em]
  \item forecast action は \fillin{\sourcepath{ros2 launch}}{22mm} ではなく \fillin{\sourcepath{fetch\_weather\_forecast.py}}{40mm} を呼ぶ。
  \item route waypoint CSV が存在しないときに使う緯度経度は
  \fillin{fallback 緯度経度}{30mm} である。
\end{enumerate}

\section{identify action と改良版 vehicle identification}
\subsection{identify action の軽量 pipeline}
\sourcepath{SolarSim.ps1 -Action identify} は、install 側の
\sourcepath{run\_identification\_pipeline.py} を呼ぶ。
これは入力ディレクトリ内のファイル有無に応じて、次を順に回す軽量 pipeline である。

\begin{enumerate}[leftmargin=1.5em]
  \item \sourcepath{gps\_track.csv} があれば route / profile 同定
  \item \sourcepath{battery\_rest.csv} があれば OCV curve 構築
  \item \sourcepath{battery\_pulse.csv} があれば Rint map 構築
  \item \sourcepath{panel\_sweep.csv} があれば panel / MPPT map 構築
  \item \sourcepath{drive\_timeseries.csv} があれば vehicle parameter fit
\end{enumerate}

\subsection{改良版 vehicle identification の主戦力}
現在の改良版では、generic replay MLE を扱う
\sourcepath{scripts/run\_vehicle\_identification.py} の方が重要である。
これは次を入口にする。
\begin{itemize}[leftmargin=1.5em]
  \item package 内 \sourcepath{data/identification/identification\_manifest.yaml}
  \item \sourcepath{normalized\_replay\_log\_csv}
  \item 既存または理論根拠のある \sourcepath{maps/*.csv}
\end{itemize}

\subsection{template と入力雛形}
\begin{center}
\begin{tabularx}{\linewidth}{>{\raggedright\arraybackslash}p{0.56\linewidth} X}
\toprule
template & 用途 \\
\midrule
\sourcepath{templates/identification\_manifest\_template.yaml} & generic replay MLE の入口 manifest。 \\
\sourcepath{templates/identification/observed\_replay\_log\_template.csv} & 時刻、距離、速度、気象、battery 観測を揃えた replay 入力。 \\
\sourcepath{templates/identification/battery\_pulse\_template.csv} & Rint 系の pulse 測定。 \\
\sourcepath{templates/identification/battery\_rest\_template.csv} & OCV 休止測定。 \\
\sourcepath{templates/identification/drive\_timeseries\_template.csv} & 駆動系時系列。 \\
\sourcepath{templates/identification/gps\_track\_template.csv} & route / profile 作成用 GPS。 \\
\sourcepath{templates/identification/panel\_sweep\_template.csv} & panel / MPPT map 作成用掃引。 \\
\sourcepath{templates/identification/source\_map\_catalog\_template.csv} & 既知 map の根拠整理。 \\
\bottomrule
\end{tabularx}
\end{center}

\subsection{問題 8: identification}
\begin{enumerate}[leftmargin=1.5em]
  \item 軽量 pipeline の入口 script は
  \fillin{\sourcepath{run\_identification\_pipeline.py}}{40mm} である。
  \item 改良版 generic replay MLE の入口 script は
  \fillin{\sourcepath{run\_vehicle\_identification.py}}{38mm} である。
  \item replay ベースで最初に揃えるべき雛形 CSV は
  \fillin{\sourcepath{observed\_replay\_log\_template.csv}}{46mm} である。
\end{enumerate}

\newpage
\section{dashboard と logger の見方}
\subsection{dashboard\_node が受けるもの}
\sourcepath{dashboard\_node} は少なくとも次を subscribe する。
\begin{itemize}[leftmargin=1.5em]
  \item planner: speed command、upper speed command、throttle command、drive mode、status、env、metrics、upper plan、lower plan
  \item vehicle: speed、SoC、battery 温度、battery 電流、battery 電圧、distance
  \item system: state、diag、MPC state、health
  \item location: \sourcepath{/sim/gps}
  \item trajectory: \sourcepath{/planner/trajectory}
\end{itemize}

つまり dashboard は、``planner が何を考えているか'' と ``実測がどうなっているか'' の両方を見せる。

\subsection{solar\_logger\_node が担うこと}
logger は、「後から再現・同定・異常解析できること」が役割である。
したがって、dashboard に見えているだけではなく、
\begin{itemize}[leftmargin=1.5em]
  \item speed command
  \item throttle command
  \item measured speed / distance
  \item battery channels
  \item system state / diag / health
\end{itemize}
を CSV として固定化する。

\subsection{問題 9: dashboard / logger}
\begin{enumerate}[leftmargin=1.5em]
  \item dashboard の HTTP API は \fillin{\sourcepath{/api/state}}{20mm} である。
  \item planner の配列 plan を dashboard に渡す topic は
  \fillin{\sourcepath{/planner/upper\_plan}}{34mm} と
  \fillin{\sourcepath{/planner/lower\_plan}}{34mm} である。
  \item 後解析用 CSV の主担当は \fillin{\sourcepath{solar\_logger\_node}}{30mm} である。
\end{enumerate}

\section{source ファイルと役割の対応表}
\begin{longtable}{p{0.31\linewidth}p{0.62\linewidth}}
\toprule
ファイル & 読むと何が分かるか \\
\midrule
\endhead
\sourcepath{SolarSim.ps1} & ユーザ操作の入口、WSL distro 解決、引数の受け渡し。 \\
\sourcepath{scripts/solar\_control.sh} & action と mode の最終分岐。 \\
\sourcepath{launch/solarcar\_sim.launch.py} & 現物として確認できる sim launch 構成。 \\
\sourcepath{scripts/solar\_sim.py} & 大会全体 simulation の処理順、出力ファイル、実行モデル。 \\
\sourcepath{mpc\_solarcar/mpc\_node.py} & planner の入出力 topic、上位 replan / 1 Hz lower、hierarchical 構成。 \\
\sourcepath{mpc\_solarcar/model.py} & 物理モデル。 \\
\sourcepath{mpc\_solarcar/dashboard\_node.py} & どの値が dashboard に乗るか。 \\
\sourcepath{mpc\_solarcar/distance\_node.py} & speed -> distance 積分。 \\
\sourcepath{mpc\_solarcar/grade\_node.py} & GPS / altitude -> grade。 \\
\sourcepath{mpc\_solarcar/solar\_logger\_node.py} & 太陽車ログ CSV の列定義。 \\
\sourcepath{mpc\_solarcar/solar\_preflight\_node.py} & solar telemetry と planner の鮮度から system state を作る。 \\
\sourcepath{scripts/fetch\_weather\_forecast.py} & 予報取得のロジック。 \\
\sourcepath{scripts/run\_identification\_pipeline.py} & 軽量同定 pipeline。 \\
\sourcepath{scripts/run\_vehicle\_identification.py} & replay MLE 同定と report 生成。 \\
\bottomrule
\end{longtable}

\section{実装監査結果と読むべき姿勢}
\flowbox{
  \textbf{重要な結論:}
  現在の package は、planner、offline sim、profile、logger / dashboard / distance / grade に加え、
  次の runtime source もワークツリー直下に揃っている。
  \begin{itemize}[leftmargin=1.5em]
    \item \sourcepath{solar\_measurement.launch.py}
    \item \sourcepath{solar\_race\_live.launch.py}
    \item \sourcepath{solar\_race\_live\_wifi.launch.py}
    \item \sourcepath{speed\_command\_bridge\_node.py}
    \item \sourcepath{telemetry\_text\_bridge\_node.py}
    \item \sourcepath{wind\_correction\_node.py}
    \item \sourcepath{solar\_state\_node.py}
    \item \sourcepath{solar\_profile.py}, \sourcepath{weather\_utils.py}
  \end{itemize}
  今チームが使うときは
  \begin{enumerate}[leftmargin=1.5em]
    \item \textbf{PowerShell 入口と profile 構造を理解する}
    \item \textbf{launch と node source で topic / parameter を追う}
    \item \textbf{source inventory の hash と配布 manifest で版を照合する}
  \end{enumerate}
  という 3 段で読むのが正しい。未確認の install 側バイナリを前提にする必要はない。
}

\section{総合問題}
\subsection*{問題 10: mode 横断追跡}
次の空欄を埋めよ。
\begin{enumerate}[leftmargin=1.5em]
  \item 大会全体の detail CSV と upper plan CSV を出す主役 script は
  \fillin{\sourcepath{scripts/solar\_sim.py}}{40mm} である。
  \item 実測車速から距離へ変換する node は \fillin{\sourcepath{distance\_node}}{28mm} である。
  \item 実走中の上位再計画と 1 Hz lower command の中核 node は \fillin{\sourcepath{mpc\_node}}{22mm} である。
  \item WiFi 受送信を束ねる設定ブロックは \fillin{\sourcepath{live.wifi\_bridge}}{34mm} である。
  \item Open-Meteo 予報取得 utility は
  \fillin{\sourcepath{fetch\_weather\_forecast.py}}{38mm} である。
  \item 改良版 replay MLE 同定は
  \fillin{\sourcepath{run\_vehicle\_identification.py}}{38mm} を読む。
\end{enumerate}

\subsection*{問題 11: フローを文章で説明せよ}
次の 2 つを、自分の言葉で説明できれば、本 package の入口理解はかなり進んでいる。
\begin{enumerate}[leftmargin=1.5em]
  \item \sourcepath{SolarSim.ps1 -Action up -Mode sim} を叩いた直後から、dashboard に値が出るまでの流れ。
  \item \sourcepath{SolarSim.ps1 -Action simulate} を叩いた直後から、summary CSV / detail CSV / upper plan CSV / HTML が保存されるまでの流れ。
\end{enumerate}

\section{最後に}
本 package を運用する上で重要なのは、``MPC の数式'' だけではない。
実際には、
\begin{itemize}[leftmargin=1.5em]
  \item どの mode を叩いたか
  \item その mode がどの profile ブロックを読むか
  \item topic がどこから入り、どこへ出るか
  \item 最後にどの CSV / dashboard / command へ着地するか
\end{itemize}
を一続きで追えることが、運用の強さになる。

本書はその全体地図である。数値を代入して全サイクルを解く問題・解答・解答用紙は、
\sourcepath{docs/complete\_flow\_workbook/} に分冊して収録する。

\end{document}
